\documentclass[10pt,a4paper]{amsart}
\usepackage[margin=19mm]{geometry}
\usepackage{amsmath,amssymb,mathtools}
\usepackage[scr=rsfs,cal=euler]{mathalfa}
\usepackage{xfrac}
\usepackage[unicode,hidelinks,bookmarks=false]{hyperref}
\newcommand{\ie}{i.e.\ }
\newcommand{\cf}{cf.\ }
\newcommand{\resp}{respectively\ }
\newcommand{\Subsection}{\subsection}
\providecommand{\nobreakdash}{\nobreak\hbox{-}}
\setlength{\emergencystretch}{2em}
\multlinegap0pt
\usepackage{bm}
\usepackage{bbm}
\usepackage{dsfont}
\usepackage{xspace}
\usepackage{cancel}
\usepackage{mathtools}
\usepackage{amsthm}
\makeatletter
\@ifundefined{theorem}{\newtheorem{theorem}{Theorem}}{}
\@ifundefined{lemma}{\newtheorem{lemma}[theorem]{Lemma}}{}
\@ifundefined{proposition}{\newtheorem{proposition}[theorem]{Proposition}}{}
\makeatother
\makeatletter
\expandafter\ifx\csname customthm\endcsname\relax
\newenvironment{customthm}[1]
{\renewcommand\theinnercustomthm{#1}\innercustomthm}
\fi
\makeatother
\title[Piecewise smooth stationary Euler flows with support in a ne]{Piecewise smooth stationary Euler flows with support in a neighborhood of a helix}
\author{Daniel Peralta-Salas, Jie Wan}
\date{Source: arXiv:2607.16141v1 (CC BY 4.0)}
\begin{document}
\maketitle

\noindent\emph{Source and licence.} Piecewise smooth stationary Euler flows with support in a neighborhood of a helix. Version arXiv:2607.16141v1 (CC BY 4.0), distributed under the Creative Commons Attribution 4.0 International licence, as recorded in the arXiv licence metadata for this exact version and shown on the abstract page https://arxiv.org/abs/2607.16141v1. Full rights notice: licenses/arxiv-2607.16141-CC-BY-4.0.txt.\par\smallskip

\noindent\textit{Editorial scope.} Counted unit: the Proposition whose source label is \texttt{prop:variation} in the article (source0.tex lines 1399--1426, source label \texttt{prop:variation}), delivered together with the article's own complete original proof of it (source0.tex lines 1427--1620, one \texttt{proof} environment of 194 lines). No mathematics is written, repaired or re-derived: statement and proof are the article's own text line for line. Required context retained uncounted: eq:phi-0 (lines 689--700); eq:phi (lines 708--716); eq:Dirichlet\_varphi (lines 740--742); prop:Dirichlet (lines 758--765); prop:first-order-dirichlet-expansion (lines 834--874); eq:prho-phi (lines 978--987); lem:poisson-calculus (lines 1324--1336). Every definition, display and label of those lines is retained unchanged. Omitted from this package and not redistributed: 1--688, 701--707, 717--739, 743--757, 766--833, 875--977, 988--1323, 1337--1398, 1621--2626. Nothing inside a retained excerpt is omitted or altered: every retained body is delivered exactly as the article's lines are written, so no omission receipt of any kind accompanies this package. Citations: the retained excerpts carry no citation command at all, so no bibliography entry of the article is transcribed and no reference is redistributed. Reference adaptation: none was needed and none was made; every reference inside the delivered unit resolves inside it, so no unresolved reference or citation is produced. Printed slips kept literal: no printed word, symbol, hypothesis or number of the counted statement or of its proof was altered. Layout and class: the article's own class is \texttt{amsart} with packages including amsmath, amssymb, amstext, mathrsfs, a4wide, graphicx, bbm, verbatim, bm, tikz, pgfplots, titletoc, fontenc, color; the wrapper is the lane's pinned \texttt{amsart} editorial wrapper at 10pt on A4, which restates the theorem-like environments the retained text needs and adds the article's own macro definitions that the retained text needs (none), with extra wrapper packages (bm, bbm, dsfont, xspace, cancel, mathtools). The delivered numbering is the unit's own. Assets: no retained span contains a figure environment, a graphics-inclusion command, a PDF-page inclusion, a table or a TikZ picture; no image file is copied into this package, the package declares no asset dependency and no image file is redistributed with it. Licence: version arXiv:2607.16141v1 of this article is distributed under the Creative Commons Attribution 4.0 International licence, as recorded in the arXiv licence metadata for this exact version and shown on the abstract page https://arxiv.org/abs/2607.16141v1. A full rights notice is delivered at licenses/arxiv-2607.16141-CC-BY-4.0.txt. Characters: every retained excerpt is pure ASCII, so the delivered engine, XeLaTeX with its default Latin Modern text font, reports no missing character.
\begin{equation}\label{eq:phi-0}
\begin{split}
&\bigl((R+\varepsilon\sigma_1x)^2+(\varepsilon\sigma_2y)^2+h^2\bigr)\bigl(h^2+(\varepsilon\sigma_2y)^2\bigr)
 \sigma_1^{-2}\partial_{xx}\phi \\
&-2\bigl((R+\varepsilon\sigma_1x)^2+(\varepsilon\sigma_2y)^2+h^2\bigr)(R+\varepsilon\sigma_1x)(\varepsilon\sigma_2y)
 \sigma_1^{-1}\sigma_2^{-1}\partial_{xy}\phi\\
&+\bigl((R+\varepsilon\sigma_1x)^2+(\varepsilon\sigma_2y)^2+h^2\bigr)\bigl(h^2+(R+\varepsilon\sigma_1x)^2\bigr)
 \sigma_2^{-2}\partial_{yy}\phi\\
&\quad=\varepsilon\bigl((R+\varepsilon\sigma_1x)^2+(\varepsilon\sigma_2y)^2+3h^2\bigr)\bigl((R+\varepsilon\sigma_1x) \sigma_1^{-1}\partial_{x}\psi+\varepsilon y\partial_{y}\psi\bigr)-2h^2F(\varepsilon^2\phi)\\
&\qquad-\bigl((R+\varepsilon\sigma_1x)^2+(\varepsilon\sigma_2y)^2+h^2\bigr) \frac{\left(F^2\right)'(\varepsilon^2\phi)}{2}+\bigl((R+\varepsilon\sigma_1x)^2+(\varepsilon\sigma_2y)^2+h^2\bigr)^2 H'(\varepsilon^2\phi).
\end{split}
\end{equation} 

\begin{equation}\label{eq:phi}
\begin{split}
\Delta\phi={}&\sigma_2^2c+\varepsilon\Bigl(-2R\sigma_1^{-1}xh^2\partial_{xx}\phi+2  R\sigma_1^{-1}\sigma_2y\partial_{xy}\phi-4R\sigma_1\sigma_2^{-1}x\partial_{yy}\phi\\
&
+(R^2+3h^2)R\sigma_1^{-1}\partial_x\phi 
+4R \sigma_1\sigma_2xc-2h^2\sqrt{F_R}\Bigr)
+\mathcal{R}(\varepsilon,\phi),
\end{split}
\end{equation}

\begin{equation}\label{eq:Dirichlet_varphi}
\phi(1+\varepsilon B(\theta),\theta) = 0 .
\end{equation}

\begin{proposition}\label{prop:Dirichlet}
	For every $B$ with $||B||_{C^{s+1}}\le1$ and for all sufficiently small $\varepsilon $, there is a unique solution
	\begin{equation*}%\label{eq:phi-eps-B-exists}
	\phi=\phi_{\varepsilon,B}\simeq \phi_0\quad\hbox{in }C^{s+1}(\Omega_{\varepsilon,B}) 
	\end{equation*}
	that satisfies Equation~\eqref{eq:phi} and the Dirichlet boundary condition~\eqref{eq:Dirichlet_varphi}.
	Furthermore, $\phi_{\varepsilon,B}<0$ in $\varOmega_{\varepsilon, B}$ and $\phi_{\varepsilon,B}$ is even if $B$ is.
\end{proposition}

 \begin{proposition}\label{prop:first-order-dirichlet-expansion}
 	For $B$ in a bounded set of $C^{s+1}(\mathbb{T})$, the solution of Proposition~\ref{prop:Dirichlet} has the expansion
\begin{equation}\label{eq:phi-expansion}
\phi_{\varepsilon,B}=A_0(\rho^2-1)+\varepsilon\left[A_1(\rho^3-\rho)\cos\theta-\frac{h^2\sqrt{F_R}}{2}(\rho^2-1)-2A_0P_{\varepsilon B}B\right]+O(\varepsilon^2),
\end{equation}
where  $A_0=\frac{\sigma_2^2c}{4}$ and
\begin{align*}\label{eq:A1}
A_1:=\frac18\Big(&-4R\sigma_1^{-1}A_0h^2
-8R\sigma_1\sigma_2^{-1}A_0
+2(R^2+3h^2)R\sigma_1^{-1}A_0 +4R \sigma_1\sigma_2c\Big).
\end{align*} 	
Moreover, its first order derivatives are given by  
\begin{equation}\label{eq:px}
\partial_x\phi_{\varepsilon,B}={} 2A_0x+\varepsilon \left[  A_1(3x^2+y^2-1)-  h^2\sqrt{F_R}x  
-2 A_0\partial_x(P_{\varepsilon B}B)\right] 
+O(\varepsilon^2),\tag{4.1a}
\end{equation}
\begin{equation}\label{eq:py}
\partial_y\phi_{\varepsilon,B}={} 2A_0y+\varepsilon \left [2A_1xy- h^2\sqrt{F_R}y  
-2 A_0\partial_y(P_{\varepsilon B}B)\right ]
+O(\varepsilon^2),\tag{4.1b}
\end{equation}
\begin{equation}\label{eq:prho}
\partial_{\rho}\phi_{\varepsilon,B}={} 2A_0\rho+\varepsilon \left[  A_1(3\rho^2-1)\cos\theta-  h^2\sqrt{F_R}\rho  
-2 A_0\partial_{\rho}(P_{\varepsilon B}B)\right] 
+O(\varepsilon^2),\tag{4.1c}
\end{equation}
%\begin{align}\label{eq:grad-phi-expansion}
%\nabla\phi_{\varepsilon,B}={}&\bigl(2A_0\rho+\varepsilon A_1(3\rho^2-1)\cos\theta-\varepsilon h^2\sqrt{F_R}\rho\bigr)e_\rho\notag\\
%& -2\varepsilon A_0\nabla(P_{\varepsilon B}B)
%+\varepsilon A_1(\rho^3-\rho)\nabla\cos\theta+O(\varepsilon^2),
%\tag{4.1d}
%\end{align}
and
\begin{equation}\label{eq:Dphi-square}
|\nabla\phi_{\varepsilon,B}|^2=4A_0^2\rho^2+4\varepsilon A_0 
\left[A_1(3\rho^3-\rho)\cos\theta-h^2\sqrt{F_R}\rho^2-2A_0\rho\partial_\rho(P_{\varepsilon B}B)\right]+O(\varepsilon^2).
\tag{4.1d}
\end{equation}
The error terms are uniform in $B$ on bounded sets.
 \end{proposition}

\begin{equation}\label{eq:prho-phi}
\begin{split}
\partial_{\rho}\phi_{\varepsilon,B}\big|_{\rho=1+\varepsilon B}&=  2A_0(1+\varepsilon B(\theta))+\varepsilon \left[  2A_1 \cos\theta-  h^2\sqrt{F_R} 
-2 A_0\partial_{\rho}(P_{\varepsilon B}B)|_{\rho=1+\varepsilon B}\right] 
+O(\varepsilon^2) \\
&=2A_0 +\varepsilon \left[ 2A_0 B(\theta)+ 2A_1 \cos\theta-  h^2\sqrt{F_R} 
-2 A_0 \Lambda_{0}B  \right] 
+O(\varepsilon^2).
\end{split}
\end{equation}

\begin{lemma} \label{lem:poisson-calculus}
Let $\textbf{B}(\theta)=\sum_{n\in\mathbb Z}B_ne^{in\theta}\in C^{s+1}(\mathbb T)$.  The Cartesian derivatives of $ P_0\textbf{B} $ up to second order are given by 
\begin{equation}\label{eq:Poisson-derivatives1}
\begin{split}
\partial_xP_0\textbf{B}&=\sum_{n\in\mathbb Z\setminus\{0\}}|n|B_n\rho^{|n|-1}e^{i(n-\text{sgn}(n))\theta},\\
\partial_yP_0\textbf{B}&=\sum_{n\in\mathbb Z\setminus\{0\}}\text{sgn}(n)i|n|B_n\rho^{|n|-1}e^{i(n-\text{sgn}(n))\theta},\\
\partial_{xx}P_0\textbf{B}&=\sum_{n\in\mathbb Z\setminus\{0,\pm1\}}|n|(|n|-1)B_n\rho^{|n|-2}e^{i(n-2\text{sgn}(n))\theta},\\
\partial_{xy}P_0\textbf{B}&=\sum_{n\in\mathbb Z\setminus\{0,\pm1\}}\text{sgn}(n)i|n|(|n|-1)B_n\rho^{|n|-2}e^{i(n-2\text{sgn}(n))\theta},\\
\partial_{yy}P_0\textbf{B}&=-\sum_{n\in\mathbb Z\setminus\{0,\pm1\}}|n|(|n|-1)B_n\rho^{|n|-2}e^{i(n-2\text{sgn}(n))\theta}.
\end{split}
\end{equation}

\end{lemma}

\begin{proposition}\label{prop:variation}
	For $ \textbf{B} \in C^{s+1}(\mathbb{T})$ and any small enough $\varepsilon$,
\begin{equation*} 
\begin{split}
\Phi_{\varepsilon,B^*,\mathbf{B}} =& -2\varepsilon A_0 P_{\varepsilon B^*} \textbf{B}+ \varepsilon^2 \Bigl[ A_0\mathcal{T} \textbf{B}
-2A_1P_{\varepsilon B^*}(\cos\theta\, \textbf{B}) +h^2\sqrt{F_R}P_{\varepsilon B^*}  \textbf{B}\\
&\hspace{8em}
 +4A_0C^*P_{\varepsilon B^*}(\cos3\theta\, \textbf{B})-2A_0t^*P_{\varepsilon B^*} \textbf{B} \Bigr]
+ O(\varepsilon^3) .
\end{split}
\end{equation*}
Moreover, its first-order derivative   satisfies
\begin{align*} 
\partial_x\Phi_{\varepsilon,B^*, \textbf{B}}
=&-2\varepsilon A_0\partial_xP_{\varepsilon B^*}  \textbf{B}
+\varepsilon^2\Big[A_0\partial_x\mathcal{T} \textbf{B}-2A_1\partial_xP_{\varepsilon B^*}(  \cos\theta\, \textbf{B}) +h^2\sqrt{F_R}\partial_xP_{\varepsilon B^*}  \textbf{B}\\
&\hspace{9em}
+4A_0C^*\partial_xP_{\varepsilon B^*}(\cos3\theta\, \textbf{B})-2A_0t^*\partial_xP_{\varepsilon B^*} \textbf{B} \Bigr]
+ O(\varepsilon^3),\notag\\
\partial_y\Phi_{\varepsilon,B^*, \textbf{B}}
=&-2\varepsilon A_0\partial_yP_{\varepsilon B^*}  \textbf{B}
+\varepsilon^2\Big[A_0\partial_y\mathcal{T} \textbf{B}-2A_1\partial_yP_{\varepsilon B^*}(  \cos\theta\, \textbf{B}) +h^2\sqrt{F_R}\partial_yP_{\varepsilon B^*}  \textbf{B}\\
&\hspace{9em}
+4A_0C^*\partial_yP_{\varepsilon B^*}(\cos3\theta\, \textbf{B})-2A_0t^*\partial_yP_{\varepsilon B^*} \textbf{B} \Bigr]
+ O(\varepsilon^3).\notag 
\end{align*}

\end{proposition} 

\begin{proof}
 

Differentiating Equation \eqref{eq:phi-0}  with respect to $ B $ at $ B = B^* $, we obtain that
$ \Phi=\Phi_{\varepsilon,B^*,\mathbf{B}} $ satisfies the equation 
 \begin{align}\label{eq:Phi}
 &\bigl((R+\varepsilon\sigma_1x)^2+(\varepsilon\sigma_2y)^2+h^2\bigr)\bigl(h^2+(\varepsilon\sigma_2y)^2\bigr)\sigma_1^{-2}\partial_{xx}\Phi \notag\\
 &  -2\bigl((R+\varepsilon\sigma_1x)^2+(\varepsilon\sigma_2y)^2+h^2\bigr)(R+\varepsilon\sigma_1x)(\varepsilon\sigma_2y)
 \sigma_1^{-1}\sigma_2^{-1}\partial_{xy}\Phi \notag\\
 & +\bigl((R+\varepsilon\sigma_1x)^2+(\varepsilon\sigma_2y)^2+h^2\bigr)\bigl(h^2+(R+\varepsilon\sigma_1x)^2\bigr)\sigma_2^{-2}\partial_{yy}\Phi \notag\\
 =&\varepsilon\bigl((R+\varepsilon\sigma_1x)^2+(\varepsilon\sigma_2y)^2+3h^2\bigr)\bigl((R+\varepsilon\sigma_1x)\sigma_1^{-1}\partial_x\Phi+\varepsilon y\partial_y\Phi\bigr)\notag\\
 &
 - 2\varepsilon^2h^2
 F'(\varepsilon^2\phi_{\varepsilon,B^*})\Phi-\varepsilon^2\bigl((R+\varepsilon\sigma_1x)^2+(\varepsilon\sigma_2y)^2+h^2\bigr)
 \left(\frac{F^2}{2}\right)''(\varepsilon^2\phi_{\varepsilon,B^*})\Phi \notag\\
 & +\varepsilon^2\bigl((R+\varepsilon\sigma_1x)^2+(\varepsilon\sigma_2y)^2+h^2\bigr)^2
 H''(\varepsilon^2\phi_{\varepsilon,B^*})\Phi.
 \end{align}
Using the fact that
\begin{equation*}\label{eq:low-regularity-FH}
F'(\varepsilon^2\phi_{\varepsilon,B^*})=O(\varepsilon),
\qquad \left(\frac{F^2}{2}\right)''(\varepsilon^2\phi_{\varepsilon,B^*})=O(1),
\qquad H''(\varepsilon^2\phi_{\varepsilon,B^*})=O(1),
\end{equation*}
Equation \eqref{eq:Phi} can be rewritten as 
\begin{equation*}\label{eq:Phi-expanded-coefficients}
\begin{split}
&\bigl(1+2R\varepsilon\sigma_1^{-1}xh^2+O(\varepsilon^2)\bigr)\partial_{xx}\Phi
-2\bigl(\sigma_2R\varepsilon\sigma_1^{-1}y+O(\varepsilon^2)\bigr)\partial_{xy}\Phi
+\bigl(1+4R\varepsilon\sigma_1\sigma_2^{-1}x+O(\varepsilon^2)\bigr)\partial_{yy}\Phi\\
&\qquad=\varepsilon\bigl(R^2+3h^2+O(\varepsilon)\bigr)\left( R\sigma_1^{-1}\partial_x\Phi
+O(\varepsilon)(\partial_x\Phi+\partial_y\Phi)\right) +O(\varepsilon^2)\Phi,
\end{split}
\end{equation*}
which implies that
\begin{align}\label{eq:Phi-expanded}
\Delta\Phi={}&\varepsilon\Big[-2R\sigma_1^{-1}h^2x\partial_{xx}\Phi
+2 R\sigma_1^{-1}\sigma_2y\partial_{xy}\Phi
-4R\sigma_1\sigma_2^{-1}x\partial_{yy}\Phi +(R^2+3h^2)R\sigma_1^{-1}\partial_x\Phi\Big] \notag\\
&
+O(\varepsilon^2)\left (\Phi+\partial_{x}\Phi+\partial_{y}\Phi+\partial_{xx}\Phi+\partial_{xy}\Phi+\partial_{yy}\Phi\right )
\end{align} 
in $ \varOmega_{\varepsilon,B^*} $. 	Likewise, differentiating the boundary condition we obtain that 
\begin{equation*} 
 \Phi(1+\varepsilon B^*(\theta),\theta)
=-\varepsilon\partial_\rho\phi_{\varepsilon,B^*}(1+\varepsilon B^*(\theta),\theta)\mathbf{B}(\theta) .
\end{equation*}
In view of the asymptotics for $ \phi_{\varepsilon,B^*}$ computed in Proposition~\ref{prop:first-order-dirichlet-expansion} and formula \eqref{eq:prho-phi}, this boundary condition can be written as 
 \begin{align*}%\label{eq:Phi-bc}
 &\Phi(1+\varepsilon B^*(\theta),\theta)\notag\\
 &\quad=-\varepsilon \textbf{B}(\theta)\bigl(2A_0+2\varepsilon A_1\cos\theta-\varepsilon h^2\sqrt{F_R}-6C^*A_0\varepsilon \cos3\theta+2A_0\varepsilon B^*(\theta)+O(\varepsilon^2)\bigr)\notag\\
 &\quad=-2\varepsilon A_0  \textbf{B}-2\varepsilon^2A_1\cos\theta\, \textbf{B}
 +\varepsilon^2h^2\sqrt{F_R} \textbf{B}+4\varepsilon^2C^*A_0\cos3\theta\, \textbf{B}-2\varepsilon^2A_0 t^*\textbf{B}+O(\varepsilon^3),
 \end{align*}
and  $ \Phi $ has the expression
 \begin{equation*}\label{eq:Phi-leading}
 \Phi=-2\varepsilon A_0P_{\varepsilon B^*} \textbf{B}+O(\varepsilon^2).
 \end{equation*}
 Equation~\eqref{eq:Phi-expanded} is then of the form
\begin{align*} 
\Delta\Phi&=\varepsilon\left [-2R\sigma_1^{-1}h^2x\partial_{xx}\Phi
+2 R\sigma_1^{-1}\sigma_2y\partial_{xy}\Phi
-4R\sigma_1\sigma_2^{-1}x\partial_{yy}\Phi +(R^2+3h^2)R\sigma_1^{-1}\partial_x\Phi\right ]  
+O(\varepsilon^3)\\
&=\varepsilon\left [I_0x\partial_{xx}\Phi
+I_1y\partial_{xy}\Phi
+I_2x\partial_{yy}\Phi +I_3\partial_x\Phi\right ]  
+O(\varepsilon^3).
\end{align*} 
% \textbf{Step 1: Low-order expansion of the variation equation.}
Let us set  $ \Phi_1:=(\Phi-\varepsilon\Phi_0)/ \varepsilon^2 $, with $ \Phi_0:=-2A_0P_{\varepsilon B^*} \textbf{B}. $ A short calculation shows that $\Phi_1$ must solve the equation
 \begin{equation}\label{eq:Phi1-equation}
 \begin{split}
 \Delta\Phi_1=&I_0x\partial_{xx}\Phi_0+I_1y\partial_{xy}\Phi_0+I_2x\partial_{yy}\Phi_0+I_3\partial_x\Phi_0+O(\varepsilon)\\
=&-2A_0I_0x\partial_{xx}P_0\textbf{B}-2A_0I_1y\partial_{xy}P_0\textbf{B}-2A_0I_2x\partial_{yy}P_0\textbf{B}-2A_0I_3\partial_xP_0\textbf{B}+O(\varepsilon) 
 \end{split}
 \end{equation}
with the boundary condition  
 \begin{equation*} 
 \Phi_1(1+\varepsilon B^*,\theta)
 =-2A_1 \cos\theta\, \textbf{B}+h^2\sqrt{F_R} \textbf{B}+4 A_0C^*\cos3\theta\, \textbf{B}-2 A_0 t^*\textbf{B}+O(\varepsilon).
 \end{equation*}
From Lemma \ref{lem:poisson-calculus}, the Fourier   terms in  Equation \eqref{eq:Phi1-equation} are obtained from
 \begin{align*} 
 x\partial_{xx}P_0\textbf{B}
 &=\sum_{n\in\mathbb{Z}\setminus\{0,\pm1\}}|n|(|n|-1)B_n\rho^{|n|-1}
 \frac{e^{i(n+1-2\text{sgn}(n))\theta}+e^{i(n-1-2\text{sgn}(n))\theta}}{2},
\notag\\
 y\partial_{xy}P_0\textbf{B}
 &=\sum_{n\in\mathbb{Z}\setminus\{0,\pm1\}}\text{sgn}(n)|n|(|n|-1)B_n\rho^{|n|-1}
 \frac{e^{i(n+1-2\text{sgn}(n))\theta}-e^{i(n-1-2\text{sgn}(n))\theta}}{2},
\notag\\
 x\partial_{yy}P_0\textbf{B}
 &=-\sum_{n\in\mathbb{Z}\setminus\{0,\pm1\}}|n|(|n|-1)B_n\rho^{|n|-1}
 \frac{e^{i(n+1-2\text{sgn}(n))\theta}+e^{i(n-1-2\text{sgn}(n))\theta}}{2},
\notag\\
 \partial_xP_0\textbf{B}
 &=\sum_{n\in\mathbb{Z}\setminus\{0\}}|n|B_n\rho^{|n|-1}e^{i(n-\text{sgn}(n))\theta},
\notag
 \end{align*}
for $ \textbf{B}=\sum_{n\in\mathbb Z}B_ne^{in\theta} $.  Taking these into \eqref{eq:Phi1-equation}, we get
\begin{equation*} 
 \Delta\Phi_1=D_0+O(\varepsilon), 
\end{equation*}
where 
\begin{equation*} 
\begin{split}
D_0:=&-2A_0(I_0-I_2)\sum_{n\in\mathbb{Z}\setminus\{0,\pm1\}}|n|(|n|-1)B_n\rho^{|n|-1}
\frac{e^{i(n+1-2\text{sgn}(n))\theta}+e^{i(n-1-2\text{sgn}(n))\theta}}{2}\\
&-2A_0I_1\sum_{n\in\mathbb{Z}\setminus\{0,\pm1\}}\text{sgn}(n)|n|(|n|-1)B_n\rho^{|n|-1}
\frac{e^{i(n+1-2\text{sgn}(n))\theta}-e^{i(n-1-2\text{sgn}(n))\theta}}{2}\\
&-2A_0I_3\sum_{n\in\mathbb{Z}\setminus\{0\}}|n|B_n\rho^{|n|-1}e^{i(n-\text{sgn}(n))\theta}.
\end{split}
\end{equation*}
%Note that for $ k,m\in \mathbb{N} $
%\begin{equation}\label{eq:laplace-polar-monomial}
%\Delta(\rho^me^{ik\theta})=(m^2-k^2)\rho^{m-2}e^{ik\theta}.
%\end{equation}
Let us note that  $ \Delta \frac{B_n}{(|n|+2)^2-m^2}\rho^{|n|+2}e^{im\theta}=B_n\rho^{|n|}e^{im\theta}  $, provided $(|n|+2)^2\ne m^2$. 
If we set  
 \begin{align}\label{eq:Phi2}
 \Phi_2:=&\sum_{|n|\ge3}\frac{(|n|-1)(I_0-I_2+I_1)+2I_3}{4}A_0B_n
 (\rho^{|n|-1}-\rho^{|n|+1})e^{i(n-\text{sgn} (n))\theta}\notag\\
 &+\sum_{|n|\ge3}\frac{|n|(I_0-I_2-I_1)}{8}A_0B_n
 (\rho^{|n|-3}-\rho^{|n|+1})e^{i(n-3\text{sgn} (n))\theta}\notag\\
 &+\frac{I_0-I_2+I_1+2I_3}{4}A_0(B_2e^{i\theta}+B_{-2}e^{-i\theta})(\rho-\rho^3)\notag\\
 &+\frac{I_0-I_2-I_1}{4}A_0(B_2e^{-i\theta}+B_{-2}e^{i\theta})(\rho-\rho^3)
 +\frac12I_3A_0(B_1+B_{-1})(1-\rho^2), 
 \end{align}
% The special modes $n=\pm1,\pm2$ are listed separately because the high-mode formula fails at these indices.   The low modes $n=\pm1,\pm2$ are separated to avoid this resonance.
% For $|n|\ge3$,
% \begin{align}\label{eq:Phi2-laplace-check}
% \Delta\{(\rho^{|n|-1}-\rho^{|n|+1})e^{i(n-\text{sign} n)\theta}\}
% &=-4|n|\rho^{|n|-1}e^{i(n-\text{sign} n)\theta},\\
% \Delta\{(\rho^{|n|-3}-\rho^{|n|+1})e^{i(n-3\text{sign} n)\theta}\}
% &=-8|n|\rho^{|n|-1}e^{i(n-3\text{sign} n)\theta}.
% \end{align}
% The special low modes are checked in the same way.  Since
% \begin{equation}\label{eq:low-mode-deltas}
% \Delta\big((\rho-\rho^3)e^{\pm i\theta}\big)=-8\rho e^{\pm i\theta},
% \qquad
% \Delta(1-\rho^2)=-4,
% \end{equation}
% the third and fourth lines of \eqref{eq:Phi2} supply the missing $n=\pm2$ pieces, while the last line supplies the $n=\pm1$ contribution.  This proves, term by term, that
then direct computation shows that $  \Delta\Phi_2=D_0. $ 
 The boundary value of $\Phi_2$ is obtained as follows. From the Taylor formula $ \rho^m\big|_{\rho=1+\varepsilon B^*}=1+m\varepsilon B^*+O(\varepsilon^2) $, we have 
 \begin{equation*} 
 (\rho^{m_1}-\rho^{m_2})\big|_{\rho=1+\varepsilon B^*}
 =(m_1-m_2)\varepsilon B^*+O(\varepsilon^2).
 \end{equation*}
 Applying this formula to each term in Equation \eqref{eq:Phi2}, we get
 \begin{align*}\label{eq:Phi2-boundary-expanded}
 \Phi_2(1+\varepsilon B^*,\theta)={}&\varepsilon A_0B^*\sum_{|n|\ge3}
 \frac{(|n|-1)(I_0-I_2+I_1)+2I_3}{2}B_n e^{i(n-\text{sgn} (n))\theta}\notag\\
 &+\varepsilon A_0B^*\sum_{|n|\ge3}
 \frac{|n|(I_0-I_2-I_1)}{2}B_n e^{i(n-3\text{sgn} (n))\theta}\notag\\
 &+\varepsilon A_0B^*\frac{I_0-I_2+I_1+2I_3}{2}(B_2e^{i\theta}+B_{-2}e^{-i\theta}) \notag\\
 &+\varepsilon A_0B^*\frac{I_0-I_2-I_1}{2}(B_2e^{-i\theta}+B_{-2}e^{i\theta})
 +\varepsilon I_3A_0B^*(B_1+B_{-1})+O(\varepsilon^2)\\
 =&O(\varepsilon),
 \end{align*}
where we have used   $ ||\textbf{B}||_{C^s(\mathbb{T})}\leq 1 $ for $ s>2 $ and Schauder estimates.
Consequently, the function $ \Phi_3:=\Phi_1-\Phi_2 $ satisfies the equation
\begin{equation*} 
\Delta\Phi_3=O(\varepsilon)\quad\hbox{in }\varOmega_{\varepsilon,B^*}
\end{equation*} 
and the boundary condition
 \begin{align*} 
 \Phi_3(1+\varepsilon B^*,\theta)=&=-2A_1 \cos\theta\, \textbf{B}+h^2\sqrt{F_R} \textbf{B}+4 A_0C^*\cos3\theta\, \textbf{B}-2 A_0 t^*\textbf{B}+O(\varepsilon).
 \end{align*}
 This shows that
 \begin{align*}%\label{eq:Phi3-Poisson}
 \Phi_3=-2A_1P_{\varepsilon B^*}(\cos\theta\, \textbf{B})
 +h^2\sqrt{F_R}P_{\varepsilon B^*}  \textbf{B}
 +4A_0C^*P_{\varepsilon B^*}(\cos3\theta\, \textbf{B})-2A_0t^*P_{\varepsilon B^*} \textbf{B}+O(\varepsilon),
 \end{align*}
and therefore
\begin{equation*} 
\begin{split}
\Phi_1=&\Phi_2+\Phi_3\\
=&-2A_1P_{\varepsilon B^*}(\cos\theta\, \textbf{B})
+h^2\sqrt{F_R}P_{\varepsilon B^*}  \textbf{B}
+4A_0C^*P_{\varepsilon B^*}(\cos3\theta\, \textbf{B})-2A_0t^*P_{\varepsilon B^*} \textbf{B}\\
&+\sum_{|n|\ge3}\frac{(|n|-1)(I_0-I_2+I_1)+2I_3}{4}A_0B_n
(\rho^{|n|-1}-\rho^{|n|+1})e^{i(n-\text{sgn} (n))\theta}\\
&+\sum_{|n|\ge3}\frac{|n|(I_0-I_2-I_1)}{8}A_0B_n
(\rho^{|n|-3}-\rho^{|n|+1})e^{i(n-3\text{sgn} (n))\theta}\\
&+\frac{I_0-I_2+I_1+2I_3}{4}A_0(B_2e^{i\theta}+B_{-2}e^{-i\theta})(\rho-\rho^3)\\
&+\frac{I_0-I_2-I_1}{4}A_0(B_2e^{-i\theta}+B_{-2}e^{i\theta})(\rho-\rho^3)
+\frac12I_3A_0(B_1+B_{-1})(1-\rho^2)+O(\varepsilon), 
\end{split}
\end{equation*} 
and $ \Phi=\varepsilon\Phi_0+\varepsilon^2\Phi_1 $. The derivative formulas are obtained by differentiating the expansion, as claimed.
\end{proof} 

\end{document}
